% Preambolo LaTeX aggiuntivo per app_50_user_guide.
%
% Serve a una cosa sola: stampare l'appendice con il log di sessione in
% ORIZZONTALE. Le righe della tabella sono larghe 118 caratteri, e su A4
% verticale (160 mm utili) starebbero solo a 6pt - illeggibili. In orizzontale
% ci sono 257 mm, quindi 9pt bastano con margine.
%
% pdflscape ruota anche la PAGINA nel PDF, non solo il contenuto: chi legge a
% schermo la trova già nel verso giusto senza doverla girare a mano.

\usepackage{pdflscape}
\usepackage{fancyvrb}

% ⚠️ pdflscape RUOTA la pagina nel PDF, ma NON allarga l'area di testo: quella
% resta larga quanto la imposta geometry per il corpo del documento (160 mm).
% Senza newgeometry la tabella continua a uscire dal foglio, semplicemente su
% una pagina girata - il che è peggio, perché sembra a posto.

% Il log è incluso da un file esterno (logs/…): così esiste in UN posto solo e
% non può divergere dalla copia nel PDF. VerbatimInput lo prende alla lettera -
% nessuna interpretazione di markup, che su un log è esattamente ciò che serve.
\newcommand{\sessionlog}[1]{%
  \clearpage
  % ⚠️ newgeometry PRIMA di \begin{landscape}, e con l'opzione "landscape":
  % dentro l'ambiente arriva troppo tardi e l'area di testo resta larga quanto
  % nel corpo del documento. Il risultato e' una pagina girata con il log
  % stretto in una colonna al centro - a prima vista sembra funzionare.
  % Su A4 girato ci sono 297 mm, 267 utili con margini da 15: a \footnotesize
  % Courier sono ~160 caratteri, contro i 118 della riga piu' larga del log.
  \newgeometry{top=15mm,bottom=15mm,left=15mm,right=15mm,landscape}
  \begin{landscape}
  \VerbatimInput[fontsize=\footnotesize,
                 fontfamily=courier,
                 numbers=none,
                 frame=none,
                 baselinestretch=0.9]{#1}
  \end{landscape}
  \restoregeometry
  \clearpage
}
